\begin{table}[htbp]
\centering
\caption{Expected calibration error and Brier score of the detector score and of the two gates. Temperature scaling uses the parameter of the benchmark table of the main article. The gate columns are five-seed means, and lower is better for both metrics.}
\label{tab:calibration_family}
\footnotesize
\setlength{\tabcolsep}{4pt}
\resizebox{\textwidth}{!}{%
\begin{tabular}{lrrrrrrrrrr}
\toprule
\multirow{2}{*}{Dataset} & \multicolumn{5}{c}{Expected calibration error} & \multicolumn{5}{c}{Brier score} \\
\cmidrule(lr){2-6}\cmidrule(lr){7-11}
 & Raw score & Temperature & Isotonic & KG gate (B4) & Full gate (B6) & Raw score & Temperature & Isotonic & KG gate (B4) & Full gate (B6) \\
\midrule
OIA-DDR test & 0.1182 & 0.0836 & 0.0367 & 0.0266 & 0.0223 & 0.2107 & 0.2039 & 0.1974 & 0.1842 & 0.1798 \\
IDRiD & 0.1108 & 0.1035 & 0.0397 & 0.0455 & 0.0422 & 0.2207 & 0.2205 & 0.2090 & 0.2037 & 0.1993 \\
TJDR & 0.0537 & 0.0622 & 0.1507 & 0.1472 & 0.1499 & 0.1036 & 0.1048 & 0.1287 & 0.1262 & 0.1258 \\
FGADR & 0.1017 & 0.1273 & 0.2508 & 0.2136 & 0.2037 & 0.1290 & 0.1313 & 0.1821 & 0.1669 & 0.1599 \\
\bottomrule
\end{tabular}}
\par\smallskip{\footnotesize The target is the RETAIN versus SUPPRESS governance label of Section 4.4 of the main article, and every fitted mapping comes from the OIA-DDR validation split. Raw and temperature-scaled scores are min-max rescaled to the unit interval, which is not a probability transform, while the gate sigmoids are read directly.}
\end{table}
